\documentclass{beamer}
\usepackage[utf8]{inputenc}
\usepackage[T1]{fontenc}
\usetheme{Madrid}

\DeclareUnicodeCharacter{202F}{\,} % espace fine = espace fine LaTeX

\title{Mettre à jour un projet avec Git et GitHub}
\author{Bastien Marguet}
\date{\today}

\begin{document}
	
	\begin{frame}
		\titlepage
	\end{frame}
	
	% -------------------------------------------------
	\begin{frame}{Pourquoi utiliser Git et GitHub ?}
		\begin{itemize}
			\item Git permet de suivre l’évolution de ton projet fichier par fichier.
			\item GitHub héberge ton dépôt en ligne pour le partager ou le publier (ex: site GitHub Pages).
			\item La combinaison rend le déploiement de modifications rapide, sécurisé et collaboratif.
		\end{itemize}
	\end{frame}
	
	% -------------------------------------------------
	\begin{frame}{1. Se placer dans le bon dossier}
		\textbf{Commande :}
		\begin{block}{}
			\texttt{cd /chemin/vers/ton/dossier}
		\end{block}
		\begin{itemize}
			\item Ouvre ton terminal (ou Git Bash).
			\item Navigue jusqu’au dossier local où se trouvent tes fichiers du projet.
			\item Vérifie que tu vois bien tes fichiers en tapant \texttt{ls} ou \texttt{dir}.
		\end{itemize}
	\end{frame}
	
	% -------------------------------------------------
	\begin{frame}{2. Vérifier les modifications}
		\textbf{Commande :}
		\begin{block}{}
			\texttt{git status}
		\end{block}
		\begin{itemize}
			\item Affiche quels fichiers sont modifiés, ajoutés ou supprimés.
			\item Vérifie avant de préparer ton commit.
			\item Important pour éviter d’oublier un fichier important.
		\end{itemize}
	\end{frame}
	
	% -------------------------------------------------
	\begin{frame}{3. Préparer les fichiers à envoyer}
		\textbf{Commande :}
		\begin{block}{}
			\texttt{git add .}
		\end{block}
		\begin{itemize}
			\item \texttt{git add .} ajoute \textbf{tous} les fichiers modifiés au prochain commit.
			\item Tu peux être plus précis : \texttt{git add index.html style.css}.
			\item Bonne pratique : vérifier après avec \texttt{git status}.
		\end{itemize}
	\end{frame}
	
	% -------------------------------------------------
	\begin{frame}{4. Créer un commit explicite}
		\textbf{Commande :}
		\begin{block}{}
			\texttt{git commit -m "Ton message clair"}
		\end{block}
		\begin{itemize}
			\item Décris ce que tu as fait : « Mise à jour responsive », « Ajout du menu burger ».
			\item Évite les messages trop vagues comme « update ».
			\item Chaque commit est une sauvegarde cliquable.
		\end{itemize}
	\end{frame}
	
	% -------------------------------------------------
	\begin{frame}{5. Envoyer en ligne (push)}
		\textbf{Commande :}
		\begin{block}{}
			\texttt{git push origin main}
		\end{block}
		\begin{itemize}
			\item \texttt{push} = envoyer les modifications vers GitHub.
			\item \texttt{origin} = le nom du dépôt distant.
			\item \texttt{main} = ta branche principale (parfois \texttt{master}).
			\item Après le push : ton site ou ton projet est mis à jour en ligne.
		\end{itemize}
	\end{frame}
	
	% -------------------------------------------------
	\begin{frame}{Conseils pour s’adapter à chaque cas}
		\begin{itemize}
			\item Toujours vérifier le nom de ta branche : \texttt{main} ou \texttt{master}.
			\item Vérifie l’URL du dépôt distant : \texttt{git remote -v}.
			\item Si plusieurs personnes travaillent dessus, utilise \texttt{git pull} pour récupérer avant de pousser.
			\item Pour les sites GitHub Pages : vérifier sur \texttt{https://username.github.io/repo}.
		\end{itemize}
	\end{frame}
	
	% -------------------------------------------------
	\begin{frame}{Résumé : cycle basique}
		\begin{enumerate}
			\item Aller dans le bon dossier.
			\item Vérifier les changements.
			\item Ajouter les fichiers.
			\item Faire un commit avec un message clair.
			\item Pousser sur GitHub.
		\end{enumerate}
		\vspace{1em}
		\textbf{Ton projet est à jour et versionné !}
	\end{frame}


% -------------------------------
\begin{frame}{`git status` : que fait-il vraiment ?}
	\begin{itemize}
		\item \textbf{`git status`} regarde uniquement l’état local du projet.
		\item Il compare tes fichiers réels avec l’historique Git (\texttt{.git/}).
		\item Pas de connexion Internet : tout se passe sur ton disque.
		\item Utile pour savoir :
		\begin{itemize}
			\item Ce qui est modifié.
			\item Ce qui est prêt à être sauvegardé (staged).
			\item Ce qui reste à ajouter.
		\end{itemize}
	\end{itemize}
\end{frame}

% -------------------------------
\begin{frame}{Où Git enregistre l’historique ?}
	\begin{itemize}
		\item Chaque projet Git contient un dossier caché \texttt{.git/}.
		\item Ce dossier stocke :
		\begin{itemize}
			\item L’historique de tous les commits.
			\item L’adresse du dépôt distant (GitHub, GitLab…).
			\item Les branches locales et distantes.
		\end{itemize}
		\item \textbf{`git status`} lit seulement ce qu’il y a ici.
	\end{itemize}
\end{frame}

% -------------------------------
\begin{frame}{`git remote` : le lien vers GitHub}
	\begin{itemize}
		\item Un \textbf{`remote`} est une adresse vers un dépôt en ligne.
		\item Par convention, le dépôt principal s’appelle \texttt{origin}.
		\item Quand tu clones :
		\begin{block}{Exemple}
			\texttt{git clone https://github.com/user/mon-projet.git}
		\end{block}
		\item Git enregistre automatiquement :
		\begin{block}{}
			\texttt{origin = https://github.com/user/mon-projet.git}
		\end{block}
	\end{itemize}
\end{frame}

% -------------------------------
\begin{frame}{Vérifier ton `remote`}
	\textbf{Commande :}
	\begin{block}{}
		\texttt{git remote -v}
	\end{block}
	
	\begin{itemize}
		\item Montre les URLs configurées pour \texttt{fetch} (récupérer) et \texttt{push} (envoyer).
		\item Exemple de sortie :
		\begin{block}{}
			\texttt{origin  https://github.com/user/mon-projet.git (fetch)}\\
			\texttt{origin  https://github.com/user/mon-projet.git (push)}
		\end{block}
		\item Vérifie toujours ton `remote` avant un `push`.
	\end{itemize}
\end{frame}

% -------------------------------
\begin{frame}{Modifier ou ajouter un `remote`}
	\begin{itemize}
		\item Pour changer l’URL :
		\begin{block}{}
			\texttt{git remote set-url origin https://github.com/...}
		\end{block}
		\item Pour ajouter un `remote` si besoin :
		\begin{block}{}
			\texttt{git remote add origin https://github.com/...}
		\end{block}
		\item Pour supprimer :
		\begin{block}{}
			\texttt{git remote remove origin}
		\end{block}
	\end{itemize}
\end{frame}

% -------------------------------
\begin{frame}{Résumé : `status` vs `remote`}
	\begin{itemize}
		\item \textbf{`git status`} = vérifie l’état local.
		\item \textbf{`remote`} = configure où envoyer et récupérer.
		\item \texttt{git push origin main} :
		\begin{itemize}
			\item Utilise le `remote` nommé \texttt{origin}.
			\item Pousse sur la branche \texttt{main} sur GitHub.
		\end{itemize}
		\item Toujours vérifier \texttt{remote -v} avant de pousser.
	\end{itemize}
\end{frame}





% -------------------------------
\begin{frame}{n8n}
	Test\_n8n\\
tjcl avho kuao znef \\

Twilio : sandbox \\
Account SID : ACa595994b97d09cc3d143d389e63c4a4f\\
Auth Token : b2ac6127ce51d80dbc40e214f8f3230d \\
\end{frame}


\end{document}
